\chapter*{第一卷既有译稿说明}
\addcontentsline{toc}{chapter}{第一卷既有译稿说明}
\markboth{译稿说明}{译稿说明}
本译稿以用户上传的 \texttt{EGAI.pdf} 为底本。原著为 A. Grothendieck 与 J. Dieudonné
合作撰写的 \textit{Éléments de géométrie algébrique：I. Le langage des schémas}，
刊于 \textit{Publications mathématiques de l'I.H.É.S.} 第 4 卷（1960），5—228 页。
本书中的未来计划、历史叙述及跨卷引文，均保留原著写作时的语境。

\medskip
\noindent\textbf{译法与记号。}
沿用“典范”“归纳极限”“射影极限”“诺特”“有限型”“拟凝聚层”“凝聚层”等译法。
层的一点处的 \textit{fibre} 译为“茎”，截面的 \textit{germe} 译为“芽”；
态射的 \textit{fibre} 译为“纤维”。
\textit{préschéma} 与 \textit{schéma} 分别译为“预概形”与“概形”，
不将本版的定义替换为其他版本的定义。
\textit{préadique} 与 \textit{adique} 分别译为“预进制”与“进制”。
条目及公式编号依照原文；不带章号的交叉引用指其所在章，“（0，…）”指第 0 章。
中文分页不与原文逐页对应；各小节注明原文印刷页码以便对照。

\medskip
\noindent\textbf{原文与译注。}
译文保留原著的组织、证明详略与引文。原著只引用、略去或留给读者的证明，不自行补写。
原作者的脚注标为“原注”；译者对术语、记号和可辨疑点的说明标为“译注”。
扫描文本的识别错误按页面图像恢复；涉及原版自身疑点的处理须另作明确说明，
不将译者的解释冒充原著内容。

\medskip
\noindent\textbf{版本说明。}
使用实际的 \texttt{ctexbook} 文档类及其默认字体配置，以 XeLaTeX 编译。
中文翻译与排版由 ChatGPT 生成，供研读及原文对照使用，未经原作者或出版机构审定。
\par\medskip\noindent\textbf{收录范围。}本合订本收入本卷导言、第 0 章 §§1—7、第 I 章 §§1—10，以及原著参考文献、记号索引、术语索引和书末资料。


\medskip
\noindent\textbf{V2 校订范围（2026 年 10 月 6 日）。}
本次在第 0 章（4.1.5）、（6.7.8），以及第 I 章（5.3.9）、（7.4.4）
修订或扩充译注：分别给出外幂、忠实平坦性、任意无挠层情形的反例，并补足对角像局部闭性的论证。
反例与补证均属译者说明，原著正文没有因此被静默改写。
另为正文中可按章号和条目号确定目标的交叉引文增加点击跳转；
跳转表示该编号的查阅位置，不代表已独立验证每个引文在论证中的适用性。
原著索引中的旧编号保留原样，不为其自动建立可能误导的链接。

本版附独立的《第一轮校订报告》，列出反例的完整计算、保留的待核问题及实际编译检查结果。
全卷结构和链接目标的核对，不等同于全卷逐句翻译复核，也不等同于所有数学证明已获独立审定。

% END frontmatter_full.tex



